% Import with \input{.../paper_threshold_figures.tex}.
% Requires \usepackage{graphicx} and \usepackage{subcaption}.
% Override this path before \input if needed.
\providecommand{\taskThresholdFigureDir}{outputs/task3/threshold/paper_figures}

\begin{figure*}[p]
  \centering
  \captionsetup{font=footnotesize}
  \captionsetup[subfigure]{font=footnotesize,skip=1pt}
  \begin{subfigure}{\textwidth}
    \centering
    \includegraphics[width=0.84\linewidth]{\taskThresholdFigureDir/paper_threshold_baselines.pdf}
    \caption{Selected organizer baselines.}
    \label{fig:task2-threshold-baselines}
  \end{subfigure}

  \begin{subfigure}{\textwidth}
    \centering
    \includegraphics[width=0.84\linewidth]{\taskThresholdFigureDir/paper_threshold_participants.pdf}
    \caption{Selected participant submissions.}
    \label{fig:task2-threshold-participants}
  \end{subfigure}

  \caption{Task~2 threshold sensitivity for selected organizer baselines (a)
  and participant submissions (b). Rows show
  \texttt{eng\_Latn--arz\_Arab}, \texttt{eng\_Latn--ekk\_Latn}, and
  \texttt{zho\_Hans--jpn\_Jpan}; columns span contrasting submitted
  precision/recall operating points. Curves give MCC
  (blue), precision (green), and recall (orange). Curves omit only the
  zero-true-positive high-threshold tail. Vertical red and blue dashed lines mark submitted and per-language optimal thresholds;
  annotations give threshold and MCC changes. Dotted gray lines in (a) give matching
  Task~1-to-Task~3 MCC. Scores use the submitted 0--100 scale and error-free is
  positive. Annotations report numeric submitted and selected threshold values
  without distinguishing inclusive and exclusive rules. Optima are selected
  and evaluated on the same assessments, so improvements are post-hoc
  optimistic upper bounds.}
  \label{fig:task2-thresholds}
\end{figure*}
